\section{Introduction}
\pagenumbering{arabic}
\setcounter{page}{1}

\subsection{Background of the Study}
Diabetes is a chronic disease characterized by elevated blood sugar levels and remains a significant public health concern in the Philippines. Individuals with diabetes face higher risks of cardiovascular disease, kidney failure, nerve damage, and vision impairment \cite{Khokhar2025}. Recent national data reveal that the prevalence of diabetes among Filipino adults has reached approximately 7.5\% as of 2024, with over 4.7 million cases reported and an alarming number of undiagnosed individuals, surpassing 53\% of the diabetes-affected population \cite{IDFPhilippines2025}. In 2020 alone, diabetes accounted for 37,265 deaths, ranking as the fourth leading cause of mortality in the country \cite{Cudis2021}. The increasing urbanization, sedentary lifestyles worsened by COVID-19 lockdowns, and changing dietary patterns have accelerated obesity and other risk factors among Filipinos, intensifying the diabetes epidemic \cite{Cando2024}.

Despite numerous international studies on diabetes risk assessment, most rely on data from Western or global populations \cite{Kiran2025}, failing to account for the unique Filipino demographic, genetic, and environmental factors \cite{Cando2024}. This reliance on international datasets undermines the applicability and reliability of predictive models for local populations, which presents a distinctive research gap. The need for a customized, locally relevant platform for accurate diabetes risk prediction is critical to improving public health outcomes \cite{Khokhar2025}. Recent advancements in artificial intelligence (AI) and machine learning have opened new avenues for medical diagnostics and predictive healthcare. However, transparency and interpretability remain major challenges for clinical adoption. Explainable Artificial Intelligence (XAI) bridges this gap by providing interpretable predictions that empower healthcare practitioners and patients alike. 

In the Philippine context, there is a pressing need for predictive models tailored to local data to ensure confidence in clinical decisions. This study addresses this gap by developing a machine learning system that integrates three key components: comprehensive evaluation of multiple machine learning models (including tree-based algorithms and deep learning approaches) trained exclusively on Philippine-specific data to identify the optimal predictor; Explainable Artificial Intelligence (XAI) techniques, specifically LIME and SHAP, to provide granular interpretations of predictions; and post-hoc calibration mechanisms to align predicted probabilities with actual clinical risk levels. By systematically combining rigorous model comparison, interpretability, and probabilistic calibration, this research aims to foster robust clinical confidence, thereby reducing undiagnosed cases and supporting effective diabetes prevention strategies within Filipino communities.


\subsection{Statement of the Problem}
Diabetes is a significant health challenge in the Philippines, and machine learning models have shown promise in predicting its risk. While Explainable Artificial Intelligence (XAI) techniques, such as LIME and SHAP, help provide transparency into model predictions, they often fail to address the issue of model calibration. Uncalibrated models produce probability estimates that may be overconfident or underconfident, leading to unreliable risk predictions. This undermines the potential of AI to support clinical decision-making and reduce diabetes incidence. Therefore, there is a critical need for calibration in XAI to ensure that predicted probabilities align with actual risks. By integrating calibration techniques into XAI models, this study aims to improve both the trustworthiness and interpretability of diabetes risk predictions. This research will contribute to developing a Calibrated Explainable AI model tailored to Philippine data, ultimately enhancing predictive accuracy and fostering greater confidence in AI-driven healthcare applications.

\subsection{Research Questions}

This study aims to answer the following research questions:

\begin{enumerate}
    \item Which among the evaluated machine learning models --- Naive Bayes, k-Nearest Neighbors, Random Forest, XGBoost, LightGBM, CatBoost, AdaBoost, Bayesian-Optimized TabNet, and ensemble methods (Voting, Stacking) --- achieves the best diabetes risk prediction performance on the DOST-FNRI 2018--2021 Expanded National Nutrition Survey dataset, as measured by F2-score and AUPRC?

    \item How do probability calibration techniques (Platt Scaling and Isotonic Regression) affect the reliability of predicted diabetes risk probabilities from the best-performing model, as assessed by Brier Score, Expected Calibration Error (ECE), and calibration plots?

    \item What are the most influential clinical, anthropometric, biochemical, and dietary features driving individual diabetes risk predictions in the Philippine context, as revealed by SHAP and LIME explanations applied to the best-performing calibrated model?

    \item How can the best-performing calibrated model with XAI explanations be effectively delivered through a web-based decision support tool that provides real-time diabetes risk assessment with interpretable, probability-calibrated outputs for end-users?
\end{enumerate}

\subsection{Objectives of the Study}

The study aims to facilitate early detection and prediction of diabetes risk among individuals in the Philippines. Calibration techniques will be applied to improve the reliability of predicted probabilities, and XAI methods will be integrated to provide interpretable explanations. The specific objectives of this study are as follows:

\begin{enumerate}
    \item Allows the primary user (healthcare practitioner or end-user) to:
    \begin{enumerate}
        \item Input personal health information relevant to diabetes risk prediction (e.g., age, gender, medical history, anthropometric and biochemical measurements).
        \item Request a prediction of the likelihood of developing diabetes based on the input data.
        \item Receive an interpretable explanation of the prediction through XAI methods, specifically SHAP (SHapley Additive exPlanations) and LIME (Local Interpretable Model-agnostic Explanations), to help understand the key factors influencing the model's output.
        \item View the model's output in the form of a calibrated diabetes risk probability and risk level classification, supporting informed decision-making and preventive actions.
    \end{enumerate}

    \item Allows the system to:
    \begin{enumerate}
        \item Perform comprehensive data pre-processing on the DOST-FNRI 2018--2021 ENNS dataset, including missing value imputation, target variable encoding, feature scaling (MinMaxScaler), correlation-based feature elimination, and class imbalance handling (SMOTE).
        \item Apply hyperparameter tuning via RandomizedSearchCV (optimizing for F2-score) to the following models: k-Nearest Neighbors, AdaBoost, XGBoost, Random Forest, LightGBM, and CatBoost; and Bayesian Optimization via Optuna for TabNet.
        \item Evaluate and compare the performance of all trained models using Precision, Recall, F2-score, AUPRC, and ROC-AUC, where F2-score serves as the primary metric.
        \item Implement and compare the following machine learning models:
        \begin{enumerate}
            \item Traditional ML: Naive Bayes, k-Nearest Neighbors, Random Forest, XGBoost, LightGBM, CatBoost, AdaBoost
            \item Deep Learning: Bayesian-Optimized TabNet (BO-TabNet)
            \item Ensemble Methods: Soft Voting Ensemble, Stacking Ensemble (with Logistic Regression meta-learner)
        \end{enumerate}
        \item Apply both Platt Scaling and Isotonic Regression to calibrate the predicted probabilities of trained models using a held-out validation set, and evaluate calibration quality using Brier Score and Expected Calibration Error (ECE).
        \item Apply SHAP and LIME to the best-performing calibrated model to generate global feature importance rankings and local per-instance explanations.
    \end{enumerate}

    \item Allows the user to:
    \begin{enumerate}
        \item Input a manual entry or upload a CSV file containing relevant health information for diabetes risk prediction.
        \item Request a diabetes risk prediction with calibrated probability estimates and XAI explanations.
        \item View the prediction results including:
        \begin{enumerate}
            \item Calibrated risk probability (e.g., ``65\% diabetes risk'')
            \item Risk level classification (Low, Moderate, or High)
            \item SHAP or LIME feature importance explanations indicating which factors most influenced the prediction
        \end{enumerate}
    \end{enumerate}
\end{enumerate}


\subsection{Significance of the Project}

This study aims to provide a \textbf{reliable and calibrated explainable AI (XAI) model} for accurate and early prediction of diabetes risk, which will significantly contribute to the following:

\begin{enumerate}
    \item \textbf{Prediabetic Patients} \\
    This study's \textit{calibrated XAI model} can provide actionable insights for prediabetic individuals, helping them make informed decisions regarding lifestyle changes to prevent or manage diabetes. The reliable predictions and interpretations provided by the model will empower patients to take timely actions toward healthier lifestyles and more proactive health management.

    \item \textbf{Medical Professionals} \\
    By integrating calibration with XAI, the study provides \textit{accurate, interpretable, and trustworthy} diabetes risk predictions that assist healthcare professionals in making informed clinical decisions. The ability to trust the \textit{probability estimates} generated by the model will facilitate personalized care, early intervention, and better patient outcomes. The explainability of the model also ensures that clinicians understand the reasoning behind the predictions, enhancing trust in the AI system.

    \item \textbf{Future Researchers in Health and Machine Learning} \\
    This study contributes to the field of healthcare AI by integrating calibration techniques into XAI models, making it a key reference for future researchers. The findings will offer insights into how calibration improves the reliability and clinical applicability of AI models, guiding future work in healthcare AI, particularly in ensuring that AI-based predictions are not only accurate but also trustworthy and interpretable.
\end{enumerate}


\subsection{Scope and Limitations}

The following are the limitations of this study:

\begin{enumerate}
    \item This study will only utilize local datasets of \textbf{Department of Science and Technology - Food and Nutrition Research Institute (DOST-FNRI)} to develop and validate a model for predicting the risk of diabetes.
    
    \item The study will focus on developing and validating a \textit{Calibrated Explainable AI (XAI)} model using machine learning techniques such as tree-based algorithms and deep learning approaches, integrated with calibration methods like Platt scaling and Isotonic Regression, to ensure accurate and reliable risk predictions.
    
    \item Limitations in the quality and quantity of available Philippine health datasets may impact the model's performance, \textit{especially when trying to represent diverse regions or populations}.
    
    \item While XAI techniques like LIME will be used to enhance the \textit{interpretability} of the model, some features of the model may still be difficult to explain comprehensively, potentially influencing how users (healthcare professionals and researchers) interpret model predictions.
    
    \item The \textit{calibrated XAI} model developed in this study is intended as a \textit{decision support tool} to assist healthcare professionals with diabetes risk prediction. It is not intended to replace or override professional medical judgment or clinical decisions.
    
    \item The trained model's predictive capacity is limited to \textit{predicting the general risk of diabetes} and does not have the ability to distinguish between \textit{type 1 diabetes} and \textit{type 2 diabetes} in individual predictions.
    
    \item Calibration techniques applied on a dataset from the \textbf{Philippine population} may not generalize well to other populations or settings. This is because the calibration process is specific to the data used to train the model, and demographic, lifestyle, or healthcare-related differences in other regions could lead to miscalibrated predictions. Further research and cross-validation across diverse populations may be necessary.
\end{enumerate}

\subsection{Assumptions}
\begin{enumerate}
    \item It is assumed that the DOST-FNRI 2018--2021 Expanded National Nutrition Survey (ENNS) dataset contains sufficient records and feature coverage across clinical, anthropometric, biochemical, and dietary components to train machine learning models capable of generalizing to the adult Filipino population.

    \item It is assumed that the selected machine learning algorithms are capable of learning meaningful patterns from the preprocessed ENNS dataset for binary diabetes risk classification, and that hyperparameter tuning via RandomizedSearchCV and Bayesian Optimization will yield models that perform better than default-parameter baselines.

    \item It is assumed that applying Platt Scaling or Isotonic Regression to a held-out validation set will produce calibrated probability estimates that are more reliable than the raw uncalibrated outputs, as indicated by improvements in Brier Score and Expected Calibration Error (ECE), even if perfect calibration is not achievable.

    \item It is assumed that SHAP and LIME will produce sufficiently interpretable local and global explanations to identify the primary features influencing individual diabetes risk predictions. It is acknowledged that LIME explanations may exhibit instability across similar instances and that neither method guarantees clinically definitive explanations.

    \item It is assumed that the calibrated probability outputs of the best-performing model can serve as meaningful decision-support information for end-users, while recognizing that clinical validation by qualified medical professionals remains necessary before any deployment in actual healthcare settings.

    \item It is assumed that the best-performing trained model, once serialized, can be loaded and served within a Python-based web application framework to provide predictions within a reasonable response time under standard hardware conditions.

\end{enumerate}